\documentclass[11pt,a4paper]{article}
\usepackage{fontspec,amsmath,amssymb,amsthm,microtype}
\setmainfont{lmroman10-regular.otf}[BoldFont=lmroman10-bold.otf,ItalicFont=lmroman10-italic.otf,BoldItalicFont=lmroman10-bolditalic.otf]
\usepackage{xcolor}
\usepackage[most]{tcolorbox}
\usepackage{mathrsfs}
\providecommand{\lt}{<}
\providecommand{\gt}{>}
\setmonofont{DejaVuSansMono.ttf}[BoldFont=DejaVuSansMono-Bold.ttf,ItalicFont=DejaVuSansMono-Oblique.ttf,BoldItalicFont=DejaVuSansMono-BoldOblique.ttf,Scale=0.86]
\newfontfamily\leanfallback{STIXTwoMath-Regular.otf}[Scale=0.86]
\definecolor{leanrule}{gray}{0.62}
\definecolor{leanback}{gray}{0.962}
\definecolor{leantitle}{gray}{0.22}
\definecolor{leanmuted}{gray}{0.38}
\newtcolorbox{leanbox}[1]{enhanced,breakable,colback=leanback,colframe=leanrule,
  boxrule=0.4pt,arc=0.6mm,left=2.4mm,right=2.4mm,top=1.8mm,bottom=1.8mm,
  before skip=0.8em,after skip=1em,fonttitle=\small\bfseries,coltitle=leantitle,
  colbacktitle=leanback,titlerule=0.3pt,title={#1}}
\newlength{\leancharwidth}
\newenvironment{leancode}{\par\smallskip\begingroup\ttfamily\footnotesize\raggedright
  \settowidth{\leancharwidth}{M}\setlength{\parskip}{0pt}\setlength{\parindent}{0pt}
  \frenchspacing}{\par\endgroup\smallskip}
\newcommand{\leanline}[2]{\hangindent=\dimexpr#1\leancharwidth+4\leancharwidth\relax
  \hangafter=1\noindent\hspace*{#1\leancharwidth}#2\par}
\newcommand{\leansource}[1]{\par{\footnotesize\color{leanmuted}#1}\par}
\newcommand{\leanlabel}[1]{\par\smallskip\noindent{\small\bfseries #1}\par}
\usepackage[margin=25mm]{geometry}
\usepackage[unicode,hidelinks]{hyperref}
\hypersetup{pdfauthor={LangevinLean Project}}
\setlength{\emergencystretch}{3em}
\setlength{\parskip}{0.4em}
\setlength{\parindent}{0pt}
\allowdisplaybreaks[2]
\raggedbottom
\title{\texorpdfstring{Actual pair laws determine the observable compression}{Actual pair laws determine the observable compression}}
\author{LangevinLean Project}
\date{}
\begin{document}\maketitle
Draft mathematical exposition; these notes may contain mistakes.


The compression of stationary observations depends only on their true two-time joint law. In particular equal pair laws give equal operators, and reversal symmetry of the pair proves self-adjointness.


\section{\texorpdfstring{From actual pair laws to operator equality}{From actual pair laws to operator equality}}


\textbf{Claim.} The evaluation pairing is the integral of the actual pointwise observable pairing under the joint state law. Thus equality of pair laws identifies the compressed operators. If that law is unchanged by coordinate swap, the operator is self-adjoint.



\textbf{Proof.} Compare the exact evaluation representatives, then change variables under the measurable pair map. Equal pair laws give equal integrals against every pair of square-integrable observables. The adjoint correlation identity and separation by inner products yield operator equality. Swapping the pair exchanges the two observable arguments and gives the symmetric pairing identity required for self-adjointness.

\begin{leanbox}{Lean: \texttt{evaluation\_\allowbreak{}inner\_\allowbreak{}pair}, \texttt{compression\_\allowbreak{}eq\_\allowbreak{}of\_\allowbreak{}pair\_\allowbreak{}law}, \texttt{compression\_\allowbreak{}isSelfAdjoint\_\allowbreak{}of\_\allowbreak{}pair\_\allowbreak{}swap}}
\leansource{\texttt{evaluation\_\allowbreak{}inner\_\allowbreak{}pair} \quad WassersteinMixing/\allowbreak{}SemigroupPair.lean, lines 19–29}
\begin{leancode}
\leanline{0}{theorem\ evaluation\_inner\_pair\ (X\ Y\ :\ Lp\ E\ 2\ Q)\ (hX\ :\ Q.map\ X\ =\ μ)\ (hY\ :\ Q.map\ Y\ =\ μ)}
\leanline{4}{(f\ g\ :\ Lp\ F\ 2\ μ)\ :}
\leanline{4}{inner\ {\leanfallback 𝕜}\ (evaluation\ ({\leanfallback 𝕜}\ :=\ {\leanfallback 𝕜})\ μ\ Q\ Y\ hY\ f)\ (evaluation\ ({\leanfallback 𝕜}\ :=\ {\leanfallback 𝕜})\ μ\ Q\ X\ hX\ g)\ =}
\leanline{6}{∫\ p,\ inner\ {\leanfallback 𝕜}\ (f\ p.2)\ (g\ p.1)\ ∂Q.map\ (fun\ ω\ ↦\ (X\ ω,\ Y\ ω))}
\end{leancode}
\leansource{\texttt{compression\_\allowbreak{}eq\_\allowbreak{}of\_\allowbreak{}pair\_\allowbreak{}law} \quad WassersteinMixing/\allowbreak{}SemigroupPair.lean, lines 34–47}
\begin{leancode}
\leanline{0}{theorem\ compression\_eq\_of\_pair\_law\ \{Ω'\ :\ Type*\}\ [MeasurableSpace\ Ω']\ (Q'\ :\ Measure\ Ω')}
\leanline{4}{(X\ Y\ :\ Lp\ E\ 2\ Q)\ (X'\ Y'\ :\ Lp\ E\ 2\ Q')}
\leanline{4}{(hX\ :\ Q.map\ X\ =\ μ)\ (hY\ :\ Q.map\ Y\ =\ μ)}
\leanline{4}{(hX'\ :\ Q'.map\ X'\ =\ μ)\ (hY'\ :\ Q'.map\ Y'\ =\ μ)}
\leanline{4}{(hpair\ :\ Q.map\ (fun\ ω\ ↦\ (X\ ω,\ Y\ ω))\ =\ Q'.map\ (fun\ ω\ ↦\ (X'\ ω,\ Y'\ ω)))\ :}
\leanline{4}{ObservableTransition.compression\ (evaluation\ ({\leanfallback 𝕜}\ :=\ {\leanfallback 𝕜})\ (F\ :=\ F)\ μ\ Q\ X\ hX)}
\leanline{6}{(evaluation\ μ\ Q\ Y\ hY)\ =}
\leanline{4}{ObservableTransition.compression\ (evaluation\ μ\ Q'\ X'\ hX')\ (evaluation\ μ\ Q'\ Y'\ hY')}
\end{leancode}
\leansource{\texttt{compression\_\allowbreak{}isSelfAdjoint\_\allowbreak{}of\_\allowbreak{}pair\_\allowbreak{}swap} \quad WassersteinMixing/\allowbreak{}SemigroupPair.lean, lines 52–73}
\begin{leancode}
\leanline{0}{theorem\ compression\_isSelfAdjoint\_of\_pair\_swap\ (X\ Y\ :\ Lp\ E\ 2\ Q)}
\leanline{4}{(hX\ :\ Q.map\ X\ =\ μ)\ (hY\ :\ Q.map\ Y\ =\ μ)}
\leanline{4}{(hswap\ :\ (Q.map\ (fun\ ω\ ↦\ (X\ ω,\ Y\ ω))).map\ Prod.swap\ =}
\leanline{6}{Q.map\ (fun\ ω\ ↦\ (X\ ω,\ Y\ ω)))\ :}
\leanline{4}{IsSelfAdjoint\ (ObservableTransition.compression}
\leanline{6}{(evaluation\ ({\leanfallback 𝕜}\ :=\ {\leanfallback 𝕜})\ (F\ :=\ F)\ μ\ Q\ X\ hX)\ (evaluation\ μ\ Q\ Y\ hY))}
\end{leancode}
\leanlabel{How the Lean proof follows the argument}
\texttt{evaluation\_\allowbreak{}inner\_\allowbreak{}pair} proves the actual pushforward integral identity. \texttt{compression\_\allowbreak{}eq\_\allowbreak{}of\_\allowbreak{}pair\_\allowbreak{}law} uses inner-product extensionality. \texttt{compression\_\allowbreak{}is\allowbreak{}Self\allowbreak{}Adjoint\_\allowbreak{}of\_\allowbreak{}pair\_\allowbreak{}swap} applies the proved pair reversal and the bounded symmetric-operator criterion.
\end{leanbox}

\end{document}
